\section{Arithmetic involution, descent and polarization}\label{sec:descent}

We construct two four-dimensional representations over $\Q_2$ whose
restrictions to $G_K$, after extending coefficients to $E$, are $V_4$.
Throughout this section $\chi=\chi_{-3}$ is the quadratic character of
$K/\Q$, and $\tau$ denotes the character of $\Q_2(-1)$.

\subsection{The rational ambient representation}

Over $\Q$, let
\[
 f:T=\PP^1_z\setminus\{0,\infty,\mu_3\}\longrightarrow U,
 \qquad \lambda=z^3,
 \qquad \mathcal A_0=\Sym^4(R^1\pi_*\Q_2).
\]
The divisor $\mu_3$ and the finite etale map $f$ are defined over $\Q$.
Set
\[
 H=H^1_!(T_{\overline\Q},f^*\mathcal A_0),\qquad
 A=a^*,\quad a(z)=\zeta_Kz,\qquad
 e_{\mathrm{nt}}=\frac{2-A-A^2}{3},\quad
 H_{\mathrm{nt}}=e_{\mathrm{nt}}H.
\]
Here $\zeta_K\in K$ is a chosen primitive cube root, and $A$ means
pullback, with the identity identification $a^*f^*\mathcal A_0=f^*\mathcal A_0$.
The operator $A$ is $\Q_2$-linear, although $a$ is defined over $K$.
If $R$ denotes the Galois action on $H$, conjugation of the deck map gives
\begin{equation}\label{eq:descent-deck-galois}
 A^3=1,\qquad R(g)AR(g)^{-1}=A^{\chi(g)}.
\end{equation}
Consequently $e_{\mathrm{nt}}$ and its image are $G_\Q$-stable.

Finite pushforward is exact, and the projection formula identifies
$f_*f^*\mathcal A_0$ with $\mathcal A_0\otimes f_*\Q_2$.
These assertions can be checked on the finite sums of geometric stalks;
see also \cite[Proposition 59.55.2, tag 03QP]{StacksFinite}.
Since $f$ is proper, the same identification applies to compactly
supported cohomology and commutes with $H_c^1\to H^1$.
Thus it also identifies their parabolic images. With the chosen
coefficient root $\zeta_E\in E$, take $\mathcal L_\rho$ to be the
$A=\zeta_E$ eigensheaf of $f_*E$, as in the preceding construction.
It follows that
\begin{equation}\label{eq:descent-eigenspaces}
 H_{\mathrm{nt}}\otimes_{\Q_2}E
   =V_{4,\rho}\oplus V_{4,\rho^{-1}},\qquad
 \dim_{\Q_2}H_{\mathrm{nt}}=8.
\end{equation}
The dimension follows from the previously proved dimension four of
each character summand. Both summands are unramified at primes of $K$
above $p\ge5$. Since $K/\Q$ is unramified at such $p$, the inertia
group of $\Q$ lies in $G_K$, so $H_{\mathrm{nt}}$ is also unramified there.

\subsection{An actual coefficient involution}

\begin{lemma}\label{lem:descent-involution}
There is a $\Q_2$-linear operator $I$ on $H_{\mathrm{nt}}$ such that
\[
 I^2=1,\qquad IR(g)=R(g)I,\qquad IAI^{-1}=A^{-1}.
\]
It is induced by an arithmetic coefficient isomorphism over inversion
$z\mapsto z^{-1}$ and respects the compact-support and boundary maps.
\end{lemma}

\begin{proof}
Put $i(\lambda)=1/\lambda$. On the finite etale double cover
$u^2=\lambda$ of $U$, the zero-preserving elliptic isomorphism
\begin{equation}\label{eq:descent-elliptic-map}
 \phi_u:\mathcal E_\lambda\longrightarrow\mathcal E_{1/\lambda},
 \qquad (x,y)\longmapsto(x/\lambda,y/u^3)
\end{equation}
is verified by
\[
 (x/\lambda)(x/\lambda-1)(x/\lambda-1/\lambda)
   =x(x-1)(x-\lambda)/\lambda^3=(y/u^3)^2.
\]
Replacing $u$ by $-u$ composes this map with $[-1]$, which acts as
$+1$ on $\Sym^4R^1\pi_*\Q_2$. The two pullbacks therefore agree
on the double-cover relation and descend to an isomorphism
\[
 \theta:i^*\mathcal A_0\longrightarrow\mathcal A_0
\]
over $\Q$; in fact all these maps are defined over $\Z[1/6]$.
For the second inversion choose $1/u$. Directly on both coordinates,
$\phi_{1/u}\circ\phi_u=\operatorname{id}$, whence
\begin{equation}\label{eq:descent-coefficient-cocycle}
 \theta\circ i^*\theta=\operatorname{id}.
\end{equation}
Choosing $-1/u$ gives the same identity on fourth symmetric powers.

Let $b(z)=1/z$. Since $fb=if$, pullback of $\theta$ gives
$\Theta:b^*f^*\mathcal A_0\to f^*\mathcal A_0$.
Define $I=\Theta_*b^*$ on cohomology. The map $b$ extends to an
automorphism of $\PP^1_z$ preserving its boundary as a set. Its
pullback and the coefficient isomorphism therefore act on $H_c^1$,
$H^1$, and the natural map between them; they act on the parabolic
image and on the boundary exact sequence as well. Equation
\eqref{eq:descent-coefficient-cocycle} gives $I^2=1$.
Both $b$ and $\Theta$ are defined over $\Q$, so $I$ commutes with $R$.
Finally $bab=a^{-1}$, while the coefficient map is pulled back from
the $\lambda$-line and the deck coefficient maps are identities.
This proves $IAI^{-1}=A^{-1}$ on coefficient systems and hence on
cohomology. It also proves that $I$ preserves $H_{\mathrm{nt}}$ and
exchanges the summands in \eqref{eq:descent-eigenspaces}.
\end{proof}

\subsection{The two descents and the quadratic twist}

\begin{theorem}\label{thm:descent-rational}
The spaces
\[
 W_4^+=H_{\mathrm{nt}}^{I=1},\qquad
 W_4^-=H_{\mathrm{nt}}^{I=-1}
\]
are continuous four-dimensional $\Q_2$-representations of $G_\Q$,
unramified outside $2,3$, and
\[
 (W_4^\pm|_{G_K})\otimes_{\Q_2}E\simeq V_4,
 \qquad W_4^-\simeq W_4^+\otimes\chi.
\]
\end{theorem}

\begin{proof}
First consider extension of the base field with coefficients $E$.
Let $V=\ker(A-\zeta_E)$ in $H_{\mathrm{nt}}\otimes E$, and put
$\delta(g)=(1-\chi(g))/2$. For $g\notin G_K$, both $R(g)$ and $I$
exchange the two eigenspaces. Hence
\begin{equation}\label{eq:descent-extensions}
 R_+(g)=I^{\delta(g)}R(g)|_V,\qquad
 R_-(g)=(-I)^{\delta(g)}R(g)|_V
\end{equation}
preserve $V$. The commutation in Lemma~\ref{lem:descent-involution}
and additivity of $\delta$ modulo two show that these maps are
representations. Their restriction to $G_K$ is the original action.

Coefficient descent is supplied by the displayed rational eigenspaces.
Indeed
\[
 \Phi_\pm:V\longrightarrow W_4^\pm\otimes E,
 \qquad v\longmapsto v\pm Iv
\]
are isomorphisms whose inverses are projection to $V$. For $g\notin G_K$,
\[
 R(g)(v\pm Iv)=\Phi_\pm(\pm IR(g)v);
\]
for $g\in G_K$ the intertwining is immediate. Thus these rational
spaces realize exactly \eqref{eq:descent-extensions} and have dimension
four. Continuity and the ramification assertion follow from $H_{\mathrm{nt}}$.

For an explicit coefficient-descent datum, let
$C=\operatorname{id}\otimes\sigma$, where
$\sigma\in\Gal(E/\Q_2)$ sends $\zeta_E$ to $\zeta_E^{-1}$.
This semilinear operator commutes with $A,I,R$ and exchanges the
two eigenspaces. On $V$ the operators $D_+=IC$, $D_-=-IC$ are
semilinear involutions commuting with the corresponding $R_\pm$.
Under $\Phi_\pm$ they are ordinary coefficient conjugation. This
checks the coefficient cocycle separately from base-field extension.

Finally put $J=A-A^2$ on $H_{\mathrm{nt}}$. Since $A+A^2=-1$ there,
\[
 J^2=-3,\qquad IJ=-JI,\qquad R(g)JR(g)^{-1}=\chi(g)J.
\]
Thus $J$ is an invertible $\Q_2$-linear map from $W_4^+$ to $W_4^-$,
intertwining the latter with $W_4^+\otimes\chi$.
\end{proof}

\begin{remark}
The signs refer to the literal map \eqref{eq:descent-elliptic-map}:
$W_4^+$ is its $+1$ eigenspace. Replacing $I$ by $-I$ gives $W_4^-$.
This constructs two particular extensions without any irreducibility
assumption or classification of all extensions.
\end{remark}

\subsection{Boundary action}

The boundary lines above $\lambda=1$ have trivial arithmetic coefficient
action. Formula \eqref{eq:descent-elliptic-map} extends regularly there,
since $\lambda$ and $u$ are units; at $u=1$ it specializes to the
identity. The choice $u=-1$ has the same effect on $\Sym^4$.
Consequently $I$ acts on the three boundary lines by the permutation
$z\mapsto z^{-1}$, with coefficient $+1$.
The nontrivial deck boundary is therefore the augmentation subspace
$B_{\mathrm{nt}}\subset\Q_2[\mu_3]$. The boundary contributions at
$0,\infty$ are killed by $e_{\mathrm{nt}}$, equivalently because
the two twisted eigensheaves have no inertia invariants there.
Writing $C_{\mathrm{nt}}=e_{\mathrm{nt}}H_c^1(T_{\overline\Q},f^*\mathcal A_0)$,
we obtain
\begin{equation}\label{eq:descent-boundary-exact}
 0\longrightarrow B_{\mathrm{nt}}\longrightarrow C_{\mathrm{nt}}
   \longrightarrow H_{\mathrm{nt}}\longrightarrow0.
\end{equation}
In $B_{\mathrm{nt}}$, the lines generated by
$2[1]-[\zeta_K]-[\zeta_K^2]$ and $[\zeta_K]-[\zeta_K^2]$
have $I$ eigenvalues $+1,-1$ and Galois characters $1,\chi$, respectively.
Thus, for $C_\pm=C_{\mathrm{nt}}^{I=\pm1}$,
\begin{equation}\label{eq:descent-boundary-traces}
 \begin{aligned}
 \Tr(\Frob_p\mid W_4^+)&=\Tr(\Frob_p\mid C_+)-1,\\
 \Tr(\Frob_p\mid W_4^-)&=\Tr(\Frob_p\mid C_-)-\chi(p)
 \end{aligned}
 \qquad(p\ge5).
\end{equation}
In particular the second boundary contribution is $+1$ at inert primes.

\subsection{Alternating pairing and reciprocity}

\begin{proposition}\label{prop:descent-pairing}
Each $W_4^\pm$ has a perfect alternating $G_\Q$-equivariant pairing
with values in $\Q_2(-5)$. Its similitude character is $\tau^5$ and
its determinant is $\tau^{10}$.
\end{proposition}

\begin{proof}
The elliptic cup pairing $\langle\ ,\ \rangle$ on $R^1\pi_*\Q_2$
is alternating, perfect, and valued in $\Q_2(-1)$. On $\mathcal A_0$ put
\[
 B(v_1v_2v_3v_4,w_1w_2w_3w_4)
   =\sum_{\sigma\in S_4}\prod_{j=1}^4\langle v_j,w_{\sigma(j)}\rangle.
\]
This is symmetric with values in $\Q_2(-4)$. In a symplectic basis
$e,f$, its monomial matrix is antidiagonal with entries
$(-1)^j j!(4-j)!\langle e,f\rangle^4$, so it is perfect over $\Q_2$.
The degree-one elliptic isomorphisms $\phi_u$ preserve cup product;
hence $\theta$ preserves $B$.

Poincare duality and the compact-support map give a perfect parabolic
pairing $\Omega:H\times H\to\Q_2(-5)$
\cite[Theorem 24.1]{Milne}; see also the support and trace formalism
of \cite{SGA4}. To see perfection on the image, the two maps
$H_c^1\to H^1$ for dual coefficient systems are adjoint up to the
degree-one cup sign, and their kernels are the annihilators of the opposite
images. The degree-one cup sign and symmetry of $B$ make $\Omega$
alternating. Both $A$ and $I$ preserve it: their coefficient maps
are isometries and their curve automorphisms have degree one, so
preserve the trace map on $H_c^2$.

The adjoint of $A$ is $A^{-1}$; thus $e_{\mathrm{nt}}$ is self-adjoint
and its image has nondegenerate restricted pairing. Since $I$ is an
involutive isometry, its two eigenspaces are orthogonal and each is
nondegenerate. The restricted forms give the asserted equivariant
pairings. Their multiplier is the character of $\Q_2(-5)$, namely
$\tau^5$. The square of a nondegenerate alternating two-form in
dimension four is a volume form, yielding $\det W_4^\pm=\tau^{10}$.
\end{proof}

\begin{corollary}\label{cor:descent-reciprocity}
For $p\ge5$, the characteristic polynomial of geometric Frobenius
on either constructed representation is
\[
 \det(X-\Frob_p\mid W_4^\pm)
   =X^4-c_{1,\pm}(p)X^3+c_{2,\pm}(p)X^2
       -p^5c_{1,\pm}(p)X+p^{10}.
\]
The analogous formula for $V_4$ at a good prime of $K$ uses its norm
$q$ in place of $p$.
\end{corollary}

\begin{proof}
The pairing identity $F^t\Omega F=p^5\Omega$ pairs eigenvalues with
their quotients $p^5/\alpha$, and the determinant is $p^{10}$.
These identities give the displayed coefficients. Restriction to
$G_K$ and scalar extension transfer the pairing to $V_4$.
\end{proof}

\begin{remark}
This argument establishes reciprocity without a purity premise. The
coefficients of the displayed polynomials lie in $\Q_2$; a rational
compatible-system assertion, Hodge numbers and bad-prime factors
require their own arguments. Good-prime purity is proved next by a
separate application of Deligne's weight theorem; the rational motivic
realization and Hodge numbers are established in Section~\ref{sec:scholl},
and the local factors in Sections~\ref{sec:local-three} and \ref{sec:local-two}.
\end{remark}

% Insert after the descent/pairing results; all referenced results then precede this.
% Standard-theorem consequence, not a new general purity theorem.
\subsection{Good-prime Frobenius purity}

\begin{proposition}[Good-prime purity]\label{prop:good-prime-purity}
For $r=2,4$, the representation $H_r$ is pure of weight $r+1$ at
every prime $\mathfrak p\nmid6$ of $K$. More explicitly, if
$q=N\mathfrak p$ and $\alpha$ is a geometric Frobenius eigenvalue
on $H_r$, then $\alpha$ is algebraic over $\Q$ and
$|\iota(\alpha)|=q^{(r+1)/2}$ for every complex embedding of its
number field.
\end{proposition}

\begin{proof}
Fix $\mathfrak p\nmid6$, and work on
$U_0=\PP^1_{\F_q}\setminus\{0,1,\infty\}$. This is a smooth
curve of finite type over $\F_q$, and the coefficient prime $2$
is different from its characteristic. The specialized Legendre
family is smooth and proper over $U_0$, so
$\mathcal V_0=R^1\pi_*E$ is lisse. At any closed point $x$ of
degree $d$, its geometric Frobenius eigenvalues are the two
roots of $X^2-a_{q^d}(x)X+q^d$; the Weil theorem for elliptic
curves gives algebraic eigenvalues of absolute value $q^{d/2}$
under every complex embedding. Hence $\mathcal V_0$ is
pointwise pure of weight one.

The cover $z^3=\lambda$ is finite etale on $U_0$. Its chosen
rank-one character summand is lisse and has finite monodromy,
with Frobenius eigenvalues in $\mu_3$ at every closed point;
it is therefore pointwise pure of weight zero. Thus
$\mathcal G_{r,0}=\Sym^r\mathcal V_0\otimes\mathcal L_{\rho,0}$
is lisse and pure of weight $r$. The claims hold after extending
coefficients to $\overline{\Q}_2$ and for every embedding used
to define weights. Finite coefficient extension and passage to
the specified direct summand preserve these algebraicity and absolute-value assertions.

Deligne's \emph{Weil II}, Corollary~3.3.6 \cite{DeligneWII}, states that for a
smooth finite-type scheme and a lisse pure sheaf of weight $w$,
the image of $H_c^i\to H^i$ is pure of weight $w+i$.
Applying this result with $i=1$ proves that the special-fibre
parabolic group is pure of weight $r+1$. Equivalently, on this
curve one can apply Theorem~3.2.3 to
$H^1(\PP^1_{\overline{\F}_q},j_*\mathcal G_{r,0})$.
The specialization isomorphism of Section~\ref{sec:geometry}
is Frobenius equivariant and identifies this image with $H_r$,
so the asserted arithmetic purity follows. Neither argument
requires irreducibility or semisimplicity of $H_r$.
\end{proof}

\begin{corollary}\label{cor:descent-good-prime-purity}
Both $W_4^+$ and $W_4^-$ are pure of weight five at every rational
prime $p\ge5$: each geometric Frobenius eigenvalue is algebraic
and all its complex conjugates have absolute value $p^{5/2}$.
\end{corollary}

\begin{proof}
By Theorem~\ref{thm:descent-rational}, the restrictions of
$W_4^\pm$ to $G_K$, after extending coefficients, are $H_4$.
At a split prime $p$ the purity just proved applies to
$\Frob_p$ itself. At an inert prime, $\Frob_p^2$ restricts to
the geometric Frobenius at the degree-two prime of $K$.
Thus, if $\beta$ is an eigenvalue of $\Frob_p$ on either
$W_4^\pm$, then $\beta^2$ is algebraic and every complex
conjugate has absolute value $(p^2)^{5/2}=p^5$.
It follows that $\beta$ is algebraic and all its complex
conjugates have absolute value $p^{5/2}$, as required.
\end{proof}

\begin{remark}
This proposition concerns good-prime Frobenius weights of the
specified $2$-adic representations. It does not by itself
identify their Hodge numbers, construct rational Scholl
projectors or a compatible system, determine bad-prime Euler
factors or conductors, or prove automorphy or analytic
continuation. The separate construction in Section~\ref{sec:scholl}
supplies the specified rational Chow projectors and Hodge numbers.
Sections~\ref{sec:local-three}--\ref{sec:conductors} independently determine
the bad-prime Euler factors and conductors.
\end{remark}


\subsection{Projector compatibility on the smooth locus}
\begin{lemma}\label{lem:open-projectors}
Let $\mathcal F=\mathcal E\times_U T$ and
$h:Z=\mathcal F^{\times_T4}\to T$. On relative cohomology set
\[
 e_{\rm odd}=\prod_{j=1}^4\frac{1-[-1]_j^*}{2},\qquad
 e_{\rm alt}=\frac1{24}\sum_{\sigma\in S_4}\operatorname{sgn}(\sigma)\sigma^*,
 \qquad e_4=e_{\rm odd}e_{\rm alt}.
\]
These are rational commuting idempotents and
\[
 e_4R^j h_*\Q_2=0\quad(j\ne4),\qquad
 e_4R^4h_*\Q_2\simeq\Sym^4(R^1\pi_*\Q_2)|_T.
\]
The diagonal map induced by $\phi_u$ on $u^2=\lambda$ commutes with
$e_4$ and descends on this summand to the literal coefficient map
$\theta$ of Lemma~\ref{lem:descent-involution}. On parabolic cohomology,
$P_\pm=e_{\rm nt}(1\pm I)/2$ are orthogonal rational idempotents
with images $W_4^\pm$ and sum $e_{\rm nt}$.
\end{lemma}
\begin{proof}
Elliptic inversion acts by $-1$ on $R^1$ and by $+1$ on $R^0,R^2$.
Relative Kunneth therefore shows that $e_{\rm odd}$ selects exactly
the fourfold tensor of $R^1$, in degree four. Permuting these four
odd-degree factors introduces the Koszul sign. Consequently the signed
permutation average becomes the ordinary symmetric tensor projector
under Kunneth. This explains the sign used in Scholl's projector
\cite[Section~1.1.2]{SchollMotives}.

The diagonal elliptic isomorphism commutes with individual inversions
and factor permutations. Changing $u$ to $-u$ composes it with
simultaneous inversion on the four factors, acting by $(-1)^4=1$
on the selected summand. Its descended map is therefore exactly the
fourth symmetric power used to define $\theta$; the cocycle follows
from $\phi_{1/u}\phi_u=1$. The deck map commutes with $e_4$.
Since $IAI^{-1}=A^{-1}$, $I$ commutes with $e_{\rm nt}$, and the
last assertions follow by multiplying the displayed idempotents.
\end{proof}

This lemma establishes the smooth-family compatibility. We now extend it
to rational Chow projectors on a specified compactification, including the
auxiliary-level translations and the parabolic boundary comparison.
